\appendixexercisegroup

\begin{enumerate}
\def\labelenumi{\arabic{enumi}.}
\item
  Let S be a set of signs, let A be a word in \(\mathbf{L}_{0}(\mathbf{S})\) , and B, C be two significant segments of A. Then either B is a segment of C, or C is a segment of B, or else B and C are disjoint.
\item
  Let S be a set of signs, let A be a significant word in \(\mathrm{L}_{0}(\mathrm{S})\) , and suppose that \(A = A'BA''\) , where B is significant. Show that, if C is a significant word, then the word \(A'CA''\) is significant. (Use Exercise 1.)
\item
  Let E be a set and let f be a mapping of \(E \times E\) into E (``internal law of composition''). Let S be a set of signs which is the disjoint union of E and the set consisting of the single element f.~Put \(n(f) = 1\) , and \(n(x) = 0\) for all \(x \in E\) .
\end{enumerate}

\begin{enumerate}
\def\labelenumi{(\alph{enumi})}
\item
  Let M be the set of significant words of \(L_{0}(S)\) . Show that there is a unique mapping v of M into E satisfying the following conditions: (1) \(v(x) = x\) for all \(x \in E\) ; (2) if A and B are two significant words, then \(v(fAB) = f(v(A), v(B))\) .
\item
  For each word \(A = (s_i)_{0 \leqslant i \leqslant n}\) in \(\mathbf{L}_0(\mathbf{S})\) let \(A^*\) be the word \((s_{i_k})\) , where the \(i_k\) are the indices \(i\) such that \(s_i \neq f\) , arranged in ascending order. Two words \(A, B\) of \(\mathbf{L}_0(\mathbf{S})\) are said to be similar if \(A^* = B^*\) . Show that if the law of composition \(f\) is associative, i.e., if
\end{enumerate}

\[
f (f (x, y), z) = f (x, f (y, z)),
\]

then \(v(A) = v(B)\) whenever \(A\) and \(B\) are similar significant words (``general theorem of associativity''). (A significant word

\[
A = (s _ {i}) _ {0 \leqslant i \leqslant 2 n}
\]

is said to be normal if \(s_{i}=f\) for \(i=0,2,4,\ldots,2n-2\) , and \(s_{i}\neq f\) for the other indices. Show that every significant word is similar to a unique normal word \(A'\) , and prove that \(v(A)=v(A')\) by induction on the length of A.)

\begin{enumerate}
\def\labelenumi{\arabic{enumi}.}
\setcounter{enumi}{3}
\tightlist
\item
  Let \(A\) be a term or a relation in a theory \(\mathfrak{G}\) . Consider the sequence of assemblies defined as follows. First of all, write \(A\) . If \(A\) consists of a single letter, stop there. If not, write next the assembly or assemblies antecedent to A (if A begins with a \(\tau\) , choose one of the antecedent assemblies arbitrarily). Then write the assembly or assemblies antecedent to those of the preceding assemblies which do not consist of a single letter; and so on.
\end{enumerate}

\begin{enumerate}
\def\labelenumi{(\alph{enumi})}
\item
  If we reverse the order of this sequence of assemblies, we get a formative construction.
\item
  Let B be a balanced segment of A such that no sign of B is linked in A to a sign outside B. Show that B is either a term or a relation (use (a) and Exercise 1).
\item
  Replace B in A by a term (resp. a relation) if B is a term (resp. a relation). Show that the assembly thus obtained is a term if A is a term, a relation if A is a relation.
\end{enumerate}

\begin{enumerate}
\def\labelenumi{\arabic{enumi}.}
\setcounter{enumi}{4}
\item
  Let A be an assembly in a theory C, let T be a term in G, and let x be a letter. If \((T|x)A\) is a term (resp. a relation), then A is a term (resp. a relation). (Use Exercise 4.)
\item
  A relation in a theory \(\mathcal{G}\) is said to be logically irreducible if it begins with a relational sign. Let \(R_{1}, R_{2}, \ldots, R_{n}\) be distinct logically irreducible relations in \(\mathcal{G}\) . A logical construction with base \(R_{1}, R_{2}, \ldots, R_{n}\) is any sequence \(A_{1}, A_{2}, \ldots, A_{p}\) of assemblies in \(\mathcal{G}\) such that each \(A_{i}\) satisfies one of the following conditions: (1) \(A_{i}\) is one of the relations \(R_{1}, R_{2}, \ldots, R_{n}\) ; (2) there exists an assembly \(A_{j}\) preceding \(A_{i}\) such that \(A_{i}\) is \(\neg A_{j}\) ; (3) there exist two assemblies \(A_{j}\) and \(A_{k}\) preceding \(A_{i}\) such that \(A_{i}\) is \(\vee A_{j}A_{k}\) .
\end{enumerate}

\begin{enumerate}
\def\labelenumi{(\alph{enumi})}
\tightlist
\item
  Show that the assemblies in a logical construction with base
\end{enumerate}

\[
R _ {1}, R _ {2}, \dots , R _ {n}
\]

are relations in \(\mathfrak{C}\) . A relation which appears in such a logical construction is said to be logically constructed from \(R_{1}, R_{2}, \ldots, R_{n}\) .

\begin{enumerate}
\def\labelenumi{(\alph{enumi})}
\setcounter{enumi}{1}
\item
  If R and S are logically constructed from \(R_{1}, R_{2}, \ldots, R_{n}\) , then so are \(\neg R, \vee RS, \Longrightarrow RS\) , ``R and \(S\)'' , \(R \Longleftrightarrow S\) .
\item
  Let R be a relation in C. Consider the sequence of relations defined as follows. First of all, write R. If R is logically irreducible, stop there. If not, write next the assembly or assemblies antecedent to R (which are well determined relations). Then write the assembly or assemblies antecedent to those of the preceding assemblies which are not logically irreducible; and so on. Let \(R_{1}, R_{2}, \ldots, R_{n}\) be the distinct logically irreducible relations which arise from this construction; they are called the logical components of R. Show that R is logically constructed from its logical components, but that if one of the relations of the sequence
\end{enumerate}

\(R_{1}, R_{2}, \ldots, R_{n}\) is omitted, then R is not logically constructed from the remaining \((n - 1)\) relations.

\begin{enumerate}
\def\labelenumi{(\alph{enumi})}
\setcounter{enumi}{3}
\tightlist
\item
  Let R be a relation and let \(R_{1}, R_{2}, \ldots, R_{n}\) be distinct logically irreducible relations such that (1) R is logically constructed from \(R_{1}, R_{2}, \ldots, R_{n}\) , (2) if we omit one relation from the sequence \(R_{1}, R_{2}, \ldots, R_{n}\) , then R is not logically constructed from the remaining relations. Show that \(R_{1}, R_{2}, \ldots, R_{n}\) are the logical components of R.
\end{enumerate}

\begin{enumerate}
\def\labelenumi{\arabic{enumi}.}
\setcounter{enumi}{6}
\tightlist
\item
  Let \(R_{1}, R_{2}, \ldots, R_{n}\) be distinct logically irreducible relations (Exercise 6) in a theory C. Let \(A_{1}, A_{2}, \ldots, A_{m}\) be a logical construction with base \(R_{1}, R_{2}, \ldots, R_{n}\) . Suppose that each \(R_{j}\) is endowed with one of the signs 0, 1. Then we give each \(A_{i}\) one of the signs 0, 1 according to the following rules: (1) if \(A_{i}\) is identical with \(R_{j}\) , give \(A_{i}\) the same sign as \(R_{j}\) ; (2) if \(A_{i}\) is identical with \(\neg A_{j}\) , where \(A_{j}\) precedes \(A_{i}\) , then give \(A_{i}\) the sign 1 (resp. 0) if \(A_{j}\) has the sign 0 (resp. 1); (3) if \(A_{i}\) is identical with \(\vee A_{j}A_{k}\) , give \(A_{i}\) the sign 1 if \(A_{j}\) and \(A_{k}\) both have the sign 1; otherwise give \(A_{i}\) the sign 0. (We are applying the following ``symbolical rule'':
\end{enumerate}

\[
\neg 0 = 1, \quad \neg 1 = 0, \quad \vee 1 1 = 1, \quad \vee 1 0 = \vee 0 1 = \vee 0 0 = 0.)
\]

\begin{enumerate}
\def\labelenumi{(\alph{enumi})}
\item
  Show that there is only one way of attributing a sign to each \(A_{i}\) in accordance with these rules.
\item
  If R is logically constructed from \(R_{1}\) , \(R_{2}\) , \(\ldots\) , \(R_{n}\) , then the sign carried by R is independent of the logical construction, with base
\end{enumerate}

\[
R _ {1}, R _ {2}, \dots , R _ {n},
\]

in which \(\pmb{R}\) appears.

\begin{enumerate}
\def\labelenumi{(\alph{enumi})}
\setcounter{enumi}{2}
\item
  If R and S are logically constructed from \(R_{1}, R_{2}, \ldots, R_{n}\) , and if the signs carried by R and \(\Longrightarrow RS\) are 0, then the sign carried by S is 0.
\item
  Suppose from now on that the only axioms of \(\mathfrak{C}\) are those provided from the schemes S1 through S4. Let \(\pmb{R}\) be a theorem in \(\mathfrak{C}\) , and let \(\pmb{R}_1, \pmb{R}_2, \ldots, \pmb{R}_n\) be its logical components. Show, however, that the signs 0 and 1 are given to \(\pmb{R}_1, \pmb{R}_2, \ldots, \pmb{R}_n\) , and the sign carried by \(\pmb{R}\) is 0. (Prove this first for the case when \(\pmb{R}\) is an axiom of \(\mathfrak{C}\) ; in the general case, consider a proof of \(\pmb{R}\) , and apply (c)) \(\mathfrak{C}\) .
\item
  Let R be a logically irreducible relation in C. Show that neither R nor ``not R'' is a theorem in C. In particular, C is non-contradictory (use (d)).
\item
  Let \(\mathbf{R}_1, \mathbf{R}_2, \ldots, \mathbf{R}_n\) be distinct, logically irreducible relations in \(\mathfrak{C}\) . Consider all relations of the form `` \(\mathbf{R}_1'\) or \(\mathbf{R}_2'\) or \ldots{} or \(\mathbf{R}_n''\) '' where, for each \(i, \mathbf{R}_i'\) is one of the two relations \(\mathbf{R}_i\) , ``not \(\mathbf{R}_i\) ''. Let \(S_1, S_2, \ldots, S_p\) be these relations. Let \(T_1, T_2, \ldots, T_q\) be the relations of the form `` \(S_{i_{1}}\) and \(S_{i_{2}}\) and \(\ldots\) and \(S_{i_{r}}\) '', where \(i_{1}, i_{2}, \ldots, i_{r}\) is any strictly increasing sequence of indices. Finally, let \(T_{0}\) be the relation `` \(R_{1}\) or (not \(R_{1}\) )'', which is a theorem in C. Show that every relation logically constructed from \(R_{1}, R_{2}, \ldots, R_{n}\) is equivalent in C to exactly one of the relations \(T_{0}, T_{1}, \ldots, T_{q}\) . (Show first that every relation \(R_{i}\) is equivalent in C to one of the relations \(T_{0}, T_{1}, \ldots, T_{q}\) . If R is logically constructed from \(R_{1}, R_{2}, \ldots, R_{n}\) , argue step by step on a logical construction with base \(R_{1}, R_{2}, \ldots, R_{n}\) which contains R. Finally, use (d) to prove uniqueness.)
\item
  Let R be a relation in C, and let \(R_{1}, R_{2}, \ldots, R_{n}\) be its logical components. Then R is a theorem in C if and only if, however the signs 0 and 1 are assigned to \(R_{1}, R_{2}, \ldots, R_{n}\) , the corresponding sign carried by R is 0.
\end{enumerate}

\begin{enumerate}
\def\labelenumi{\arabic{enumi}.}
\setcounter{enumi}{7}
\tightlist
\item
  Let \(R_{1}, R_{2}, \ldots, R_{n}\) be the logically irreducible relations in a theory G (Exercise 6). Assign to each of the relations \(R_{i}\) one of the signs 0, 1, 2. To each relation in a logical construction with base \(R_{1}, R_{2}, \ldots, R_{n}\) we assign one of the signs 0, 1, 2 in accordance with the following symbolical rules (Exercise 7):
\end{enumerate}

\[
\begin{array}{c} \neg 0 = 1, \quad \neg 1 = 0, \quad \neg 2 = 2; \\ \lor 0 0 = \lor 0 1 = \lor 0 2 = \lor 1 0 = \lor 2 0 = \lor 2 2 = 0; \\ \lor 1 1 = 1, \quad \lor 1 2 = \lor 2 1 = 2. \end{array}
\]

\begin{enumerate}
\def\labelenumi{(\alph{enumi})}
\item
  If R is logically constructed from \(R_{1}, R_{2}, \ldots, R_{n}\) , the sign carried by R is independent of the logical construction, with base \(R_{1}, R_{2}, \ldots, R_{n}\) , in which R appears.
\item
  Suppose that the only axioms of \(\mathfrak{G}\) are those provided by the schemes S2, S3, S4. Let \(R\) be a theorem in \(\mathfrak{G}\) . Show that however the signs 0, 1, 2 are assigned to the logical components of \(R\) , the corresponding sign carried by \(R\) is 0. On the other hand, if \(S\) is logically irreducible and has sign 2, then the sign carried by \((S\) or \(S) \Longrightarrow S\) is 2. Deduce that \(\mathfrak{G}\) is not equivalent to the theory which has the same signs as \(\mathfrak{G}\) and whose axioms are provided by the schemes S1, S2, S3, S4.
\item
  Establish an analogous result for theories based only on the schemes S1, S3, S4, or on the schemes S1, S2, S4. (Use respectively the following rules :
\end{enumerate}

\[
\begin{array}{c} \neg 0 = 1, \quad \neg 1 = 0, \quad \neg 2 = 2, \quad \lor 0 0 = \lor 0 1 = \lor 1 0 = \lor 0 2 = \lor 2 0 = 0, \\ \quad \lor 1 1 = 1, \qquad \lor 1 2 = \lor 2 1 = 1, \qquad \lor 2 2 = 1; \\ \quad \neg 0 = 1, \qquad \neg 1 = 2, \qquad \neg 2 = 0, \\ \quad \lor 0 0 = \lor 0 1 = \lor 1 0 = \lor 0 2 = \lor 2 0 = \lor 2 1 = 0, \\ \quad \lor 1 1 = \lor 1 2 = 1, \qquad \lor 2 2 = 2.) \end{array}
\]

\begin{enumerate}
\def\labelenumi{(\alph{enumi})}
\setcounter{enumi}{3}
\tightlist
\item
  Establish an analogous result for a theory based only on the schemes S1, S2, S3. (Use four signs 0, 1, 2, 3 and the following rules:
\end{enumerate}

\[
\lnot 0 = 1, \quad \lnot 1 = 0, \quad \lnot 2 = 3, \quad \lnot 3 = 0,
\]

\[
\vee 0 0 = \vee 0 1 = \vee 1 0 = \vee 0 2 = \vee 2 0 = \vee 0 3 = \vee 3 0 = \vee 2 3 = \vee 3 2 = 0,
\]

\[
\vee 1 1 = 1, \quad \vee 1 2 = \vee 2 1 = \vee 2 2 = 2, \quad \vee 1 3 = \vee 3 1 = \vee 3 3 = 3.)
\]
